% DINO-WM PointMaze on a Colab budget — engineering and results report.
%
% Figures: reports/figures/{cem_success,cem_state_dist,cem_embed_div,success_rates}.pdf
% are generated from the recorded planner logs (exact data). Two additional
% training-curve figures are produced on Drive by notebook 07; copy them into
% reports/figures/ to have them appear (the report compiles either way):
%   <Drive>/wm_poc/logs/dino_wm/_summary/pointmaze_training_curves.png
%   <Drive>/wm_poc/logs/dino_wm/_summary/pointmaze_scratch_vs_finetune.png
%
% Build: tectonic dino_wm_pointmaze_report.tex   (or pdflatex, twice)

\documentclass[11pt]{article}
\usepackage[margin=2.6cm]{geometry}
\usepackage{graphicx}
\usepackage{booktabs}
\usepackage{amsmath}
\usepackage{xcolor}
\usepackage[colorlinks=true, linkcolor=blue!50!black, urlcolor=blue!50!black, citecolor=blue!50!black]{hyperref}
\graphicspath{{figures/}}

\newcommand{\drivefigure}[3]{%
  \IfFileExists{figures/#1}{\includegraphics[width=#2]{#1}}{%
    \fbox{\parbox[c][4.2cm][c]{0.9\linewidth}{\centering\itshape
      Figure pending: copy \texttt{#3} from the Drive summary directory into
      \texttt{\detokenize{reports/figures/#1}} and rebuild.}}}}

\title{Training and Evaluating DINO-WM on PointMaze under a Colab GPU Budget:\\
       Engineering Record and Results}
\author{Thomas Georges\thanks{Engineering and experiment automation developed
        with Claude (Anthropic) in the \texttt{World\_Models\_LAS} repository.}}
\date{June 13, 2026}

\begin{document}
\maketitle

\begin{abstract}
We reproduce the predictor-training and planning pipeline of DINO-WM
\cite{dinowm} on the PointMaze task under a strict Colab budget (target:
a full experiment in 2--3 hours on an A100, with T4 fallback). The stock
pipeline required more than 30 hours per epoch on Colab hardware; we diagnose
the cause as per-sample I/O amplification compounded by online re-encoding of
frozen DINOv2 features, and eliminate it with a precomputed-latent training
path that reproduces the upstream data distribution exactly. Epoch time on a
T4 drops from a projected $\sim$20 hours to 31 minutes ($\sim$21$\times$
per-step speedup), making the full 2{,}000-rollout run, a low-data transfer
study, and CEM planning evaluations feasible on free-tier hardware. Final
results (200 planning evaluations per checkpoint): fine-tuning a low-data
(300-rollout) predictor from the full-data checkpoint reaches the converged
validation loss immediately ($\sim$30$\times$ fewer updates than scratch
training) and delivers $2.75\times$ the planning success of the scratch
baseline ($0.275$ vs.\ $0.10$, two-proportion $z{=}4.5$, $p<10^{-5}$), while
\emph{matching} the full-data model trained on $6.7\times$ more data ($0.275$
vs.\ $0.23$; difference not significant, $p{=}0.30$).
\end{abstract}

\section{Objective and starting point}

The goal of this track is a faithful, budget-feasible reproduction of the
DINO-WM training and planning loop on PointMaze: train the latent-dynamics
predictor over frozen DINOv2 patch features on 2{,}000 rollouts, study
low-data transfer (scratch vs.\ fine-tune on 300 rollouts), and evaluate
closed-loop CEM planning --- all orchestrated from a single Colab notebook
against an unmodified upstream checkout,\footnote{\texttt{gaoyuezhou/dino\_wm},
kept pristine on Drive; all behavioural changes are applied as marker-guarded,
idempotent, reversible textual patches at launch time
(Section~\ref{sec:patches}).} with every artifact persisted to Google Drive.

The starting point was a viability probe that exposed the central obstacle:
with the stock pipeline, \emph{one epoch} of full-data training took more
than 30 hours on a T4 (and would not have been acceptable on an A100 either),
independent of batch size --- 36{,}815 steps/epoch at batch~4 and 4{,}602
steps/epoch at batch~32 both projected $\sim$19--21\,h, pinning the cost to
the per-sample data path rather than per-step compute.

\section{Diagnosis: where the 30-hour epochs came from}
\label{sec:diagnosis}

Profiling the upstream data path identified three compounding causes.

\paragraph{Per-sample episode reload.} Upstream
\texttt{TrajSlicerDataset.\_\_getitem\_\_} fetches a 4-frame training window
by calling the base dataset's \texttt{\_\_getitem\_\_}, which
\texttt{torch.load}s the \emph{entire} episode image tensor
(\texttt{obses/episode\_XXX.pth}, $\sim$93 frames) and applies the
resize/normalize transform to \emph{every} frame --- then discards all but
four. At $\sim$147k windows per epoch this re-reads and re-transforms roughly
1{,}600$\times$ the dataset size per epoch, through the Google Drive FUSE
mount.

\paragraph{Online re-encoding of frozen features.} DINO-WM trains with
\texttt{train\_encoder=false}: the DINOv2 encoder is frozen, so its outputs
are constants of the dataset. The stock loop nevertheless re-encodes every
frame of every window on every epoch --- pure recomputation.

\paragraph{Operational confounders.} The OOM-era configuration had pinned
\texttt{env.num\_workers=0} (single-threaded loading), and all reads went
through Drive FUSE rather than local disk.

Batch size cannot address any of these, which the probe confirmed empirically.

\section{The latent-cache training path}
\label{sec:latent}

The fix follows from the frozen-encoder observation: \emph{encode once, reuse
forever}. Three components, installed into the upstream checkout at launch
time:

\begin{enumerate}
\item \textbf{One-time precompute}
      (\texttt{wm\_poc\_precompute\_latents.py}): each episode is read once,
      preprocessed exactly as upstream does (\texttt{default\_transform} at
      $224$px, then the \texttt{VWorldModel} resize to $196$px, yielding
      $196$ patch tokens of dimension 384), encoded under autocast through
      the frozen \texttt{dinov2\_vits14}, and written as a per-episode
      \texttt{fp16} \texttt{.npy} plus a manifest recording encoder identity,
      patch geometry, episode lengths, and dataset coverage. The cache is
      $\sim$30\,GB for all 2{,}000 episodes, lives on session-local disk
      (Drive-hosted latents would reintroduce the FUSE bottleneck; the
      scripts warn if so configured), and the precompute is idempotent and
      resumable: existing episodes are skipped, so interrupted runs continue
      and smaller runs reuse larger caches. Measured rate on a T4 (including
      the one-time Drive read of raw episodes): $\sim$0.37 episodes/s.
\item \textbf{A drop-in latent dataset}
      (\texttt{wm\_poc\_latent\_dataset.py}), selected purely by Hydra
      override (\texttt{env.dataset.\_target\_=\ldots}). It subclasses the
      upstream \texttt{PointMazeDataset}, so actions, states,
      proprioception, normalization statistics, and the seed-42 train/val
      split are \emph{bit-identical} to the online path (verified directly
      against the upstream loader on synthetic data). Its slicer mirrors
      \texttt{TrajSlicerDataset} but memory-maps only the $\sim$0.6\,MB of
      latents each window needs, instead of materializing a $\sim$13\,MB
      episode; it also resolves \texttt{TrajSubset} wrappers at construction
      time (their \texttt{\_\_getattr\_\_} recurses if pickled into
      spawn-context dataloader workers --- a latent upstream bug). The image
      transform is retained as an attribute because the upstream planner
      builds its observation preprocessor from it.
\item \textbf{An input-dispatch bypass} in
      \texttt{VWorldModel.encode\_obs}: 4-D inputs
      $(B,T,\text{patches},\text{dim})$ are recognized as precomputed
      latents and skip the encoder; 5-D image inputs take the original path.
      The real DINOv2 weights remain in the model and its checkpoints, so
      planning and any image-based evaluation are untouched.
\end{enumerate}

A guard in the command builder rejects latent training combined with the
upstream decoder at dry-run time: the decoder needs pixel targets, which the
latent dataset does not serve (this guard was added after the low-data runs
crashed in exactly that configuration).

\subsection{Throughput effect}

\begin{table}[h]
\centering
\small
\begin{tabular}{lllll}
\toprule
Configuration & Steps/epoch & Step time (T4) & Epoch (T4) & Epoch (A100, est.) \\
\midrule
Stock pipeline, batch 4              & 36{,}815 & $\sim$3\,s    & $>$30\,h          & --- \\
Stock pipeline, batch 32             & 4{,}602  & 15.9--17\,s   & 19--21\,h (proj.) & --- \\
Latent cache, batch 32, stride 1     & 4{,}602  & $\sim$0.76\,s & $\sim$58\,min     & 6--12\,min \\
Latent cache, batch 32, stride 2     & 2{,}330  & 0.72--0.76\,s & 31\,min (meas.)   & 3--6\,min \\
\bottomrule
\end{tabular}
\caption{Training throughput before and after the latent-cache path. The
remaining step cost is the depth-6 ViT predictor itself ($\sim$2 TFLOPs
forward+backward at batch 32 over $3\times196$ tokens), i.e.\ the quantity
the experiment is supposed to measure. ``Stride 2'' subsamples window start
positions (consecutive stride-1 windows overlap by 19/20 frames); it was used
for the full run on T4 only.}
\label{tab:throughput}
\end{table}

\section{Upstream patches and what each one fixed}
\label{sec:patches}

All upstream modifications are marker-guarded textual patches with
timestamped backups, applied (idempotently) by the launch scripts and covered
by source-level unit tests, including composition tests on the real upstream
source.

\begin{table}[h]
\centering
\small
\begin{tabular}{p{4.2cm}p{9.6cm}}
\toprule
Patch & Problem it fixes \\
\midrule
\texttt{mixed\_precision} & Exposes \texttt{DINO\_MIXED\_PRECISION}
  (bf16 on A100/L4, fp16 on T4 --- Turing only emulates bf16, slowly)
  through the \texttt{Accelerator}. \\
\texttt{step\_checkpointing} & Opt-in rolling intra-epoch state-dict
  checkpoint (\texttt{model\_latest\_step.pth}: weights, optimizers, RNG,
  batch index) every $N$ steps with deterministic mid-epoch resume; bounds a
  disconnection to minutes of lost work. \\
\texttt{val\_no\_grad} & Upstream validation runs the forward pass with
  autograd enabled and never backpropagates; two full batch-32 graphs remain
  live simultaneously, which OOMs a 16\,GB T4 at the first epoch boundary
  (observed: a 338\,MiB attention-matrix allocation failing with 14.4\,GiB
  already held). Wrapping validation in \texttt{torch.no\_grad()} removes the
  graphs and speeds validation up. Memory is flat across epochs with this in
  place. \\
\texttt{finetune\_init} & Upstream has no fine-tuning support at all (the
  wrapper's overrides previously had no consumer). Adds a hook at the end of
  \texttt{init\_models} that, for a fresh run with
  \texttt{++finetuning.enabled=true}, loads predictor / action-encoder /
  proprio-encoder weights from \texttt{finetuning.init\_from}, honoring
  strictness, epoch-reset, and per-component flags; resume always wins. \\
\texttt{latent\_cache} & Installs the two latent-path modules and the
  \texttt{encode\_obs} bypass of Section~\ref{sec:latent}. \\
\texttt{evaluator\_video} & Upstream records planning videos only when the
  model has a decoder (the layout includes a decoded ``imagination'' panel),
  so decoder-free checkpoints produced no videos at all. Adds a no-decoder
  branch recording the executed rollout against the goal image, with the
  imagination panel honestly blacked out. \\
\bottomrule
\end{tabular}
\caption{Upstream patch inventory.}
\label{tab:patches}
\end{table}

\section{Reliability engineering for paid GPU time}
\label{sec:reliability}

Several failures observed during real Colab sessions led to operational
hardening, since lost A100 minutes are the scarcest resource:

\begin{itemize}
\item \textbf{Checkpoint cadence.} Full epoch checkpoints every epoch, plus
      the rolling step checkpoint every 2{,}000 steps (A100) or 500 steps
      (T4): a crash or disconnection costs at most a few minutes; relaunching
      the same cell resumes deterministically (the window permutation is
      seeded and RNG state is restored).
\item \textbf{Orphaned trainers.} A double interrupt during the SIGTERM grace
      period could orphan a trainer (wandb traps SIGTERM), which then held
      $\sim$13.7\,GB of T4 memory against subsequent launches while writing
      ghost progress bars into the run logs. Interrupts now escalate to
      SIGKILL reliably, every launch first kills leftover DINO-WM workers,
      and stale logs are rotated aside with timestamps (all generations are
      preserved and merged when training curves are parsed).
\item \textbf{Self-gating launch cells.} Every launch cell checks durable
      Drive state before spending GPU time: completed training is skipped
      (highest \texttt{model\_<N>.pth} vs.\ the config's epoch target),
      partial training resumes, and planning re-triggers only when the
      recorded evaluation does not match the config's current
      \texttt{n\_evals}. A top-to-bottom notebook run is therefore
      idempotent and performs only outstanding work.
\item \textbf{Smoke crash test.} A 20-rollout configuration exercises the
      \emph{exact} production pipeline (patches, precompute, latent
      training, checkpointing, and a 2-evaluation CEM planning pass) in
      $\sim$10--15 minutes. It caught, before they could cost full-run time:
      a planner/latent-dataset incompatibility, a checkpoint-resolution bug
      in the planning command builder, and a blank-checkpoint parsing crash.
\end{itemize}

\section{Experimental protocol}

\begin{table}[h]
\centering
\small
\begin{tabular}{lllll}
\toprule
Run & Data & Epochs & Init & Planning evals \\
\midrule
Full no-decoder      & 2{,}000 + 200 rollouts & 10 & random          & 200 \\
Low-data scratch     & 300 + 60 rollouts      & 30 & random          & 200 \\
Low-data fine-tune   & 300 + 60 rollouts      & 20 & full checkpoint & 200 \\
\bottomrule
\end{tabular}
\caption{Runs entering the results. All train the predictor only
(no decoder) on cached latents with identical seeds, normalization, and
train/val splits; planning is CEM with 10 optimization iterations, 100
samples, goal horizon 5, \texttt{goal\_source=random\_state}, seed 0. The
full run was trained on a T4 with the stride-2 window subsampling of
Table~\ref{tab:throughput}; the low-data pair uses the full stride-1 window
set.}
\label{tab:runs}
\end{table}

Fine-tuning loads the predictor and action/proprio encoders from the full
run's final checkpoint and re-trains them on the low-data subset at reduced
learning rates ($10^{-5}$ predictor, $10^{-4}$ action encoder, vs.\
$5\cdot10^{-5}$ / $5\cdot10^{-4}$ from scratch); the frozen DINOv2 encoder is
shared by construction.

\section{Results}
\label{sec:results}

All numbers in this section are read directly from the recorded artifacts ---
specifically, the \emph{latest} \texttt{final\_eval} row of each run's planner
\texttt{logs.json}, at a planning sample size of 200 episodes per checkpoint.

\subsection{Training loss: fine-tuning starts where scratch ends}

\begin{figure}[h]
\centering
\drivefigure{pointmaze_training_curves.png}{0.78\linewidth}{\_summary/pointmaze\_training\_curves.png}
\caption{Training (dashed) and validation (solid) latent-prediction loss per
epoch, parsed from the per-epoch loss lines in the run logs (all rotated log
generations merged). The low-data scratch run descends from $\approx 2.3$ to
$\approx 0.02$ over 30 epochs; the fine-tune starts at $\approx 0.015$ at
epoch 1 --- already at the scratch run's converged level --- and stays there
for all 20 epochs.}
\label{fig:curves}
\end{figure}

Two observations. First, the efficiency claim: the fine-tuned predictor
reaches at epoch~1 the validation loss the scratch baseline needs
$\sim$30 epochs to attain --- pre-training delivers the converged predictor
for $\sim$30$\times$ fewer updates on the low-data task. Second, the honest
caveat: scratch \emph{does} catch up in validation loss by epoch 30 on this
data budget. The latent-prediction loss alone therefore does not separate the
two models --- the separation appears downstream, in planning.

\begin{figure}[h]
\centering
\drivefigure{pointmaze_scratch_vs_finetune.png}{0.7\linewidth}{\_summary/pointmaze\_scratch\_vs\_finetune.png}
\caption{Final validation h-step loss, scratch (blue) vs.\ fine-tune
(orange): $\approx 0.021$ vs.\ $\approx 0.012$ --- the same order of magnitude,
in sharp contrast to their $2.75\times$ gap in planning performance
(Figure~\ref{fig:success}).}
\label{fig:scratchft}
\end{figure}

\subsection{Planning success}

\begin{figure}[h]
\centering
\includegraphics[width=0.72\linewidth]{success_rates.pdf}
\caption{CEM planning success after the final optimization iteration, 200
evaluation episodes per checkpoint, with 95\% Wilson intervals. Exact values
from the planner logs: fine-tune 0.275, full-data 0.23, scratch 0.10.}
\label{fig:success}
\end{figure}

With $n{=}200$ the intervals tighten to $\pm$0.06 and two clean results
emerge. \emph{Fine-tune $>$ scratch} is decisive (0.275 vs.\ 0.10;
two-proportion $z{=}4.5$, $p<10^{-5}$): pre-training on the full corpus and
then adapting on 300 rollouts recovers $2.75\times$ the planning success of
training on those same 300 rollouts from scratch. \emph{Fine-tune vs.\
full-data} (0.275 vs.\ 0.23) is statistically indistinguishable ($p{=}0.30$)
--- fine-tuning matches a model trained on $6.7\times$ more data. (The
larger-sample pass also corrected all three point estimates downward from the
noisier 50-eval interim --- $0.34/0.26/0.12 \to 0.275/0.23/0.10$ --- but left
the ranking and both conclusions intact.) The most interesting finding is the
dissociation already noted: scratch closes most of the validation-loss gap
(0.021 vs.\ fine-tune's 0.012; Figure~\ref{fig:scratchft}) yet still trails by
$2.75\times$ in planning success --- latent-space MSE is a weak proxy for
downstream control quality, which argues for always carrying the planning
evaluation, not only the loss curves.

\subsection{Planner diagnostics: CEM is genuinely optimizing}

\begin{figure}[h]
\centering
\begin{minipage}{0.49\linewidth}
  \includegraphics[width=\linewidth]{cem_success.pdf}
\end{minipage}\hfill
\begin{minipage}{0.49\linewidth}
  \includegraphics[width=\linewidth]{cem_state_dist.pdf}
\end{minipage}
\caption{Planner-internal telemetry per CEM iteration (exact values from
\texttt{logs.json}; for relaunched evaluations the latest attempt per
iteration is shown). Left: success rate of the current best plans rises
monotonically-ish across the 10 iterations (full: $0.135 \to 0.23$; fine-tune:
$0.135 \to 0.275$; scratch: $0.03 \to 0.10$). Right: mean state distance to
goal falls correspondingly (e.g.\ fine-tune $1.45 \to 0.93$).}
\label{fig:cem}
\end{figure}

\begin{figure}[h]
\centering
\includegraphics[width=0.6\linewidth]{cem_embed_div.pdf}
\caption{Divergence between imagined and executed visual embeddings per CEM
iteration. The fine-tune and full models keep imagination closest to
reality at the end of optimization; all models improve as CEM concentrates
on better action sequences.}
\label{fig:embdiv}
\end{figure}

The iteration curves confirm the planner is exploiting the learned dynamics
rather than stumbling onto goals: success increases and goal distance
decreases monotonically over CEM iterations for every trained model, with
the scratch model improving the least --- consistent with its weaker
dynamics. (The 2-evaluation smoke models stay at zero and are excluded.)

\subsection{Planning videos}

Each planning evaluation records per-episode MP4s
(\texttt{output\_final\_<idx>\_<success|failure>.mp4}; 26 frames at 12\,fps
--- the $5\times5{+}1$-step horizon). A historical note for interpreting older
artifacts: until the \texttt{evaluator\_video} patch
(Table~\ref{tab:patches}), upstream recorded videos only for decoder-bearing
models, so the only videos on Drive predating it come from the 1-epoch
online smoke model and are \emph{not} representative of the science
checkpoints. Videos for the real checkpoints are (re)generated by every
planning evaluation from the patch onward.

\section{Threats to validity and known caveats}
\label{sec:caveats}

\begin{itemize}
\item \textbf{Sample size.} 200 planning episodes per checkpoint ($\pm0.06$
      Wilson half-width). Both headline comparisons clear the noise; the
      fine-tune-vs-full gap is genuinely within sampling error rather than
      merely under-powered.
\item \textbf{Single seed.} All runs are seed 0. The transfer claim would be
      strengthened by at least one more seed of the low-data pair.
\item \textbf{Full-run stride.} The full-data model was trained on a T4 with
      stride-2 window subsampling (half the gradient updates per epoch of
      the paper-exact stride-1 set). Its 0.23 success is therefore a lower
      bound on what an A100 stride-1 training would deliver; comparisons
      against it should be made with this in mind. (Subsequent project
      policy: GPU-tier configurations may differ only in throughput knobs,
      never in the experiment definition.)
\item \textbf{Precision.} T4 runs use fp16, A100 runs bf16 (and TF32
      matmuls); results across tiers are statistically equivalent, not
      bit-identical.
\item \textbf{Loss-curve history.} The full run's early-epoch loss lines were
      destroyed by a (since-fixed) log-rotation flaw during repeated
      relaunches; its training curve may begin after epoch 1. The model and
      all planning results are unaffected.
\item \textbf{Validation loss naming.} The summarized
      ``final\_val\_loss\_pred\_hstep'' is upstream's epoch-level validation
      loss for the no-decoder objective (latent h-step prediction MSE).
\end{itemize}

\section{Automated final pass and hand-off}

The driver notebook is now fully self-gating, so the pending ``final
statistics'' pass is a single top-to-bottom run on the next available A100
($\approx$2--2.5\,h): it skips the completed full training, tops the scratch
checkpoint's evaluation up to $n{=}200$, retrains the fine-tune freshly
(force-restart scoped to that cell) ending in a 200-evaluation planning pass,
and evaluates the full checkpoint at $n{=}200$ --- shrinking the success-rate
intervals to $\approx\pm0.06$ and regenerating per-episode videos for all
checkpoints. On a T4 the same pass is correct but slow ($\sim$10--12\,h,
with the unresumable planning blocks as the practical risk).

The trained full-data checkpoint simultaneously serves as the frozen
\emph{global model} of the local/global track (notebook 03), which loads it
in-process through the upstream loader and trains a differentiable local
surrogate against the same latent cache. The selection logic prefers
full-data checkpoints explicitly, so re-running the transfer experiments
cannot silently change the global model.

\appendix

\section{Reproduction guide}

\begin{enumerate}
\item Open \texttt{notebooks/02\_dino\_wm\_foundation.ipynb} on Colab
      (A100 preferred; T4 works with longer wall times). Run top to bottom.
\item The preflight cell selects precision and profile by GPU, pins the
      latent cache to local disk, arms stale-process cleanup and the
      checkpoint cadence, and reports free disk (the full cache needs
      $\sim$30\,GB).
\item The smoke cell crash-tests the full pipeline in $\sim$10--15 minutes;
      everything after it is self-gating and performs only outstanding work.
\item Results: \texttt{notebooks/07\_dino\_wm\_results.ipynb} (CPU runtime
      suffices) renders the summary table, training curves, success rates,
      CEM diagnostics, and embedded planning videos, and writes the figure
      files referenced by this report into the Drive summary directory.
\end{enumerate}

\section{Artifact inventory}

\begin{itemize}
\item \textbf{Wrapper library} \texttt{src/wm\_poc/dino\_wm/}: command
      builders, config validation, patch modules
      (Table~\ref{tab:patches}), run-state gates, metrics parsers
      (rotated-log merging), visualization.
\item \textbf{Launch scripts} \texttt{scripts/dino\_wm/}: experiment /
      train / fine-tune / plan / precompute runners, patch appliers, run
      summarizers; all guarded by \texttt{RUN\_DINO\_WM=1}.
\item \textbf{Configs} \texttt{configs/dino\_wm/}: smoke (latent and legacy
      online), OOM diagnostic, full no-decoder (A100 stride-1 and T4
      stride-2 profiles), low-data scratch/fine-tune (200-eval planning),
      planner ablation.
\item \textbf{Tests}: $\sim$95 unit tests over the wrapper (patch
      application and composition against real upstream source, command
      rendering, run-state gates, metrics extraction), runnable without
      torch; functional parity of the latent dataset against the upstream
      loader validated separately with torch.
\item \textbf{Checkpoints and logs} (Drive):
      \texttt{wm\_poc/checkpoints/dino\_wm/outputs/<run>/} and
      \texttt{wm\_poc/logs/dino\_wm/<run>/} including planner
      \texttt{logs.json}, per-episode videos, and the \texttt{\_summary}
      tables and figures.
\end{itemize}

\begin{thebibliography}{1}
\bibitem{dinowm} G.~Zhou, H.~Pan, Y.~LeCun, and L.~Pinto.
\emph{DINO-WM: World Models on Pre-trained Visual Features enable Zero-shot
Planning}. arXiv:2411.04983, 2024.
\end{thebibliography}

\end{document}
